\chapter{The main theorems and their reduction to a half-plane}\label{chapter-main-theorems}

Alpöge proves the Riemann Hypothesis and the Generalized Riemann Hypothesis for Dirichlet
$L$-functions, for the functions $\zeta(s)$ and $L(s,\chi)$ of Mathlib \cite{Mat20}. This chapter
states the theorems and deduces all of them from one statement: $L(s,\chi)\neq0$ for every
Dirichlet character $\chi$ and every $s$ with $\operatorname{Re}s>1/2$, except at the pole of a
trivial character. The deduction uses the functional equation of a primitive $L$-function and the
Euler factors that separate an imprimitive character from the primitive character inducing it. The
rest of the book proves the half-plane statement for primitive characters.

\section{Dirichlet characters and their $L$-functions}\label{section-mt-l-functions}

\begin{definition}\label{mt-dirichlet-character}
Let $N\geq1$ be an integer. A \emph{Dirichlet character modulo $N$} is a function
$\chi:\mathbf Z/N\mathbf Z\to\mathbf C$ such that $\chi(1)=1$, $\chi(xy)=\chi(x)\chi(y)$ for all
$x,y\in\mathbf Z/N\mathbf Z$, and $\chi(x)=0$ for every $x$ that is not a unit. For $n\in\mathbf Z$
we write $\chi(n)$ for $\chi(n\bmod N)$.
\begin{itemize}
\item The \emph{trivial character} $\mathbf 1_N$ modulo $N$ takes the value $1$ at every unit. The
  ring $\mathbf Z/1\mathbf Z$ is zero and its element $0$ is a unit, so $\mathbf 1_1$ is the only
  Dirichlet character modulo $1$, and $\mathbf 1_1(n)=1$ for every $n\in\mathbf Z$.
\item The \emph{conjugate} of $\chi$ is the Dirichlet character $\bar\chi$ modulo $N$ with
  $\bar\chi(x)=\overline{\chi(x)}$.
\item The character $\chi$ is \emph{even} if $\chi(-1)=1$ and \emph{odd} if $\chi(-1)=-1$. Every
  Dirichlet character is even or odd. We put $a_\chi=0$ if $\chi$ is even and $a_\chi=1$ if $\chi$
  is odd.
\item Let $M\geq1$ be a divisor of $N$ and let $\psi$ be a Dirichlet character modulo $M$. The
  character modulo $N$ \emph{induced} by $\psi$ is the Dirichlet character $\chi$ modulo $N$ with
  $\chi(n)=\psi(n)$ for every $n\in\mathbf Z$ prime to $N$.
\item The \emph{conductor} of $\chi$ is the least divisor $M\geq1$ of $N$ such that $\chi$ is
  induced by a Dirichlet character modulo $M$. That character is unique; we write it $\chi^*$. The
  character $\chi$ is \emph{primitive} if its conductor is $N$.
\end{itemize}
The character $\chi^*$ is primitive, $\bar\chi$ has the same conductor as $\chi$, and
$\mathbf 1_N$ has conductor $1$ \cite[Chapter 5]{Dav00}.
\end{definition}

\begin{definition}\label{mt-l-function}
Let $N\geq1$ and let $\chi$ be a Dirichlet character modulo $N$
(Definition~\ref{mt-dirichlet-character}). The \emph{$L$-function} of $\chi$ is the function
$s\mapsto L(s,\chi)$ from $\mathbf C$ to $\mathbf C$ such that:
\begin{itemize}
\item $L(s,\chi)=\sum_{n\geq1}\chi(n)n^{-s}$ for $\operatorname{Re}s>1$;
\item if $\chi\neq\mathbf 1_N$, it is holomorphic on $\mathbf C$;
\item if $\chi=\mathbf 1_N$, it is holomorphic on $\mathbf C\setminus\{1\}$ with a simple pole at
  $s=1$, and $L(1,\mathbf 1_N)$ is the conventional value of Remark~\ref{mt-value-at-one}.
\end{itemize}
By analytic continuation, these conditions determine $L(s,\chi)$ at every $s$ other than the
pole. The \emph{Riemann zeta function} is $\zeta(s)=L(s,\mathbf 1_1)$. Thus
$\zeta(s)=\sum_{n\geq1}n^{-s}$ for $\operatorname{Re}s>1$, and $\zeta(1)=(\gamma-\log4\pi)/2$,
where $\gamma$ is Euler's constant.
\end{definition}

\begin{remark}\label{mt-value-at-one}
Mathlib constructs $L(s,\chi)$ for every $s\in\mathbf C$ from Hurwitz zeta functions, and the
construction assigns a value to $L(1,\mathbf 1_N)$, where the $L$-function has its pole. For $N=1$
the value is $\zeta(1)=(\gamma-\log4\pi)/2\approx-0.977$, the constant term at $s=1$ of the
Laurent expansion of $\pi^{-s/2}\Gamma(s/2)\zeta(s)$. It is not zero, and this matters in two places.
The hypothesis $s\neq1$ of Theorem~\ref{mt-riemann-hypothesis} excludes no zero of
$\zeta$; it is part of Mathlib's statement of the Riemann Hypothesis. The hypothesis of
Proposition~\ref{mt-imprimitive} concerns every primitive character, $\mathbf 1_1$ among them, and
every $s$ with $\operatorname{Re}s>1/2$, so it includes the assertion $\zeta(1)\neq0$, which holds
for this value. For $N>1$ no statement of this chapter depends on the value of
$L(1,\mathbf 1_N)$: each one assumes $\chi\neq\mathbf 1_N$ or $s\neq1$, or concerns
$0<\operatorname{Re}s<1$.
\end{remark}

\begin{lemma}\label{mt-euler-factors}
Let $M$ and $N$ be positive integers with $M\mid N$, let $\psi$ be a Dirichlet character modulo
$M$, and let $\chi$ be the character modulo $N$ induced by $\psi$
(Definition~\ref{mt-dirichlet-character}). Let $s\in\mathbf C$, and suppose that
$\psi\neq\mathbf 1_M$ or $s\neq1$. Then
\[
 L(s,\chi)=L(s,\psi)\prod_{p\mid N}\bigl(1-\psi(p)p^{-s}\bigr),
\]
where the product runs over the primes $p$ dividing $N$. For $\operatorname{Re}s>1$ this is the
comparison of Euler products in \cite[Chapter 5]{Dav00}, and both sides are holomorphic away from
$s=1$.
\end{lemma}

\begin{remark}\label{mt-imprimitive-zeros}
For an imprimitive character $\chi$ modulo $N$, the function $L(s,\chi)$ of
Definition~\ref{mt-l-function} is the $L$-function of $\chi$ as a character modulo $N$: its
Dirichlet series and its Euler product omit every prime dividing $N$. Lemma~\ref{mt-euler-factors}
with $\psi=\chi^*$ compares it with $L(s,\chi^*)$. Let $p$ be a prime that divides $N$ but not
the conductor of $\chi$, and write $\chi^*(p)=e^{i\theta}$. The factor $1-\chi^*(p)p^{-s}$ vanishes at
the points $s=(i\theta+2\pi ik)/\log p$ with $k\in\mathbf Z$, which lie on the line
$\operatorname{Re}s=0$. So $L(s,\chi)$ has infinitely many zeros on that line whenever such a prime exists. For example,
$L(s,\mathbf 1_N)=\zeta(s)\prod_{p\mid N}(1-p^{-s})$ for $s\neq1$, and it vanishes at $2\pi ik/\log p$
for every prime $p\mid N$ and every $k\in\mathbf Z$.
\end{remark}

\begin{remark}\label{mt-trivial-zeros}
The trivial zeros of $\zeta$ are $-2,-4,-6,\ldots$, and $\zeta(0)=-1/2$ \cite[Chapter 2]{Tit86}.
Theorem~\ref{mt-riemann-hypothesis} excludes this set, written $\{-2(n+1):n\geq0\}$ as in Mathlib's
statement of the Riemann Hypothesis, and the point $s=1$ of Remark~\ref{mt-value-at-one}; since
$\zeta(0)\neq0$, the point $0$ needs no exclusion. Let $\chi\neq\mathbf 1_1$ be a Dirichlet character
modulo $N$ and $a=a_\chi$. Then $L(s,\chi)$ vanishes at $s=-a-2k$ for every integer $k\geq0$, by
\cite[Chapter 9]{Dav00} for primitive $\chi$ and by Lemma~\ref{mt-euler-factors} in general; these
are the trivial zeros of $L(s,\chi)$. When $\chi$ is primitive,
the functional equation and the nonvanishing of $L(s,\bar\chi)$ on $\operatorname{Re}s\geq1$ show
that the trivial zeros are the only zeros with $\operatorname{Re}s\leq0$
\cite[Chapters 9 and 14]{Dav00}. So for a primitive character, Theorem~\ref{mt-strip-primitive} is
the Generalized Riemann Hypothesis in its usual form. For an imprimitive character the zeros of
Remark~\ref{mt-imprimitive-zeros} also lie on $\operatorname{Re}s=0$, and the theorems for all
characters are stated on the strip $0<\operatorname{Re}s<1$ and the half-plane
$\operatorname{Re}s>1/2$.
\end{remark}

\section{The main theorems}\label{section-mt-main-theorems}

\begin{theorem}\label{mt-riemann-hypothesis}
Let $s\in\mathbf C$ with $\zeta(s)=0$, where $\zeta$ is the Riemann zeta function of
Definition~\ref{mt-l-function}. Suppose that $s\neq-2(n+1)$ for every integer $n\geq0$, and that
$s\neq1$. Then $\operatorname{Re}s=1/2$.
\end{theorem}
\leanref{Main.lean}{riemannHypothesis}
\leanref{Main.lean}{riemannZeta_eq_zero_re}
\leanref{Statements.lean}{RiemannZetaZerosOnLine}
\begin{proof}
By Theorem~\ref{mt-half-plane} for $N=1$, $\zeta(w)=L(w,\mathbf 1_1)$ is nonzero for every
$w\in\mathbf C$ with $\operatorname{Re}w>1/2$ and $w\neq1$. Proposition~\ref{mt-riemann-of-half-plane}
gives $\operatorname{Re}s=1/2$.
\end{proof}

\begin{theorem}\label{mt-half-plane}
Let $N\geq1$, let $\chi$ be a Dirichlet character modulo $N$
(Definition~\ref{mt-dirichlet-character}), and let $s\in\mathbf C$ with $\operatorname{Re}s>1/2$.
Suppose that $\chi\neq\mathbf 1_N$ or $s\neq1$. Then $L(s,\chi)\neq0$, where $L(s,\chi)$ is the
$L$-function of Definition~\ref{mt-l-function}.
\end{theorem}
\leanref{Main.lean}{LFunction_ne_zero_of_half_lt_re}
\leanref{Statements.lean}{DirichletLNonvanishingRightOfHalf}
\leanref{Development/AutomorphicArithmetic/CharacterBoxes/MainConclusion.lean}{UpperTheory.FinalTheorems.everyCharacter}
\begin{proof}
This is Theorem~\ref{asm-every-character} of Chapter~\ref{chapter-assembly}, whose hypothesis on
Galois representations holds by Theorem~\ref{gal-galois-representations-any-rc} of
Chapter~\ref{chapter-galois-representations}. In that deduction,
Theorem~\ref{asm-zero-free-primitive} treats primitive characters, and
Theorem~\ref{asm-every-character-green} passes to all characters by
Proposition~\ref{mt-imprimitive}.
\end{proof}

\begin{theorem}\label{mt-strip}
Let $N\geq1$, let $\chi$ be a Dirichlet character modulo $N$
(Definition~\ref{mt-dirichlet-character}), and let $s\in\mathbf C$ with $0<\operatorname{Re}s<1$ and
$L(s,\chi)=0$ (Definition~\ref{mt-l-function}). Then $\operatorname{Re}s=1/2$.
\end{theorem}
\leanref{Main.lean}{LFunction_eq_zero_re}
\leanref{Statements.lean}{DirichletLZerosOnLine}
\begin{proof}
Theorem~\ref{mt-half-plane} is the hypothesis of Proposition~\ref{mt-strip-of-half-plane}, which
gives $\operatorname{Re}s=1/2$.
\end{proof}

\begin{theorem}\label{mt-strip-primitive}
Let $N\geq1$, let $\chi$ be a primitive Dirichlet character modulo $N$
(Definition~\ref{mt-dirichlet-character}), and let $s\in\mathbf C$ with $0<\operatorname{Re}s<1$
and $L(s,\chi)=0$ (Definition~\ref{mt-l-function}). Then $\operatorname{Re}s=1/2$.
\end{theorem}
\leanref{Main.lean}{LFunction_eq_zero_re_of_isPrimitive}
\leanref{Statements.lean}{PrimitiveDirichletLZerosOnLine}
\begin{proof}
Theorem~\ref{mt-half-plane} is the hypothesis of
Proposition~\ref{mt-strip-primitive-of-half-plane}, which gives $\operatorname{Re}s=1/2$.
\end{proof}

\section{Reduction to primitive characters}\label{section-mt-imprimitive}

\begin{lemma}\label{mt-euler-factor-nonzero}
Let $M\geq1$, let $\psi$ be a Dirichlet character modulo $M$
(Definition~\ref{mt-dirichlet-character}), let $p$ be a prime, and let $s\in\mathbf C$ with
$\operatorname{Re}s>0$. Then $1-\psi(p)p^{-s}\neq0$.
\end{lemma}
\leanref{Development/AutomorphicArithmetic/CharacterBoxes/ClaimStrip.lean}{Claim.factor_ne_zero}
\begin{proof}
The values of $\psi$ have absolute value at most $1$, and $|p^{-s}|=p^{-\operatorname{Re}s}<1$.
So $|\psi(p)p^{-s}|<1$.
\end{proof}

\begin{proposition}\label{mt-imprimitive}
Suppose that $L(w,\psi)\neq0$ for every integer $M\geq1$, every primitive Dirichlet character
$\psi$ modulo $M$ and every $w\in\mathbf C$ with $\operatorname{Re}w>1/2$
(Definitions~\ref{mt-dirichlet-character} and~\ref{mt-l-function}). Let $N\geq1$, let $\chi$ be a
Dirichlet character modulo $N$, and let $s\in\mathbf C$ with $\operatorname{Re}s>1/2$. Suppose that
$\chi\neq\mathbf 1_N$ or $s\neq1$. Then $L(s,\chi)\neq0$.
\end{proposition}
\leanref{Development/AutomorphicArithmetic/CharacterBoxes/Imprimitive.lean}{UpperTheory.Top.imprimitive}
\begin{proof}
Let $M$ be the conductor of $\chi$ and let $\psi=\chi^*$, a primitive character modulo $M$ that
induces $\chi$. By hypothesis $L(s,\psi)\neq0$. If $\chi\neq\mathbf 1_N$ then $\psi\neq\mathbf 1_M$,
since the character induced by a trivial character is trivial. So $\psi\neq\mathbf 1_M$ or
$s\neq1$, and Lemma~\ref{mt-euler-factors} gives
\[
 L(s,\chi)=L(s,\psi)\prod_{p\mid N}\bigl(1-\psi(p)p^{-s}\bigr).
\]
Each factor of the product is nonzero by Lemma~\ref{mt-euler-factor-nonzero}, since
$\operatorname{Re}s>0$. Hence $L(s,\chi)\neq0$.
\end{proof}

\section{Reflection by the functional equation}\label{section-mt-functional-equation}

\begin{proposition}\label{mt-functional-equation}
Let $\Gamma_{\mathbf R}(s)=\pi^{-s/2}\Gamma(s/2)$. Let $N\geq1$, let $\chi$ be a primitive Dirichlet
character modulo $N$, and let $a=a_\chi$ (Definition~\ref{mt-dirichlet-character}). The function
$\Lambda(s,\chi)=\Gamma_{\mathbf R}(s+a)L(s,\chi)$ extends to a holomorphic function on
$\mathbf C\setminus\{0,1\}$, and to an entire function if $N>1$. For every
$s\in\mathbf C\setminus\{0,1\}$,
\[
 \Lambda(1-s,\chi)=N^{s-1/2}\,\varepsilon(\chi)\,\Lambda(s,\bar\chi),\qquad
 \varepsilon(\chi)=\frac{\tau(\chi)}{i^{a}N^{1/2}},\qquad
 \tau(\chi)=\sum_{x\in\mathbf Z/N\mathbf Z}\chi(x)e^{2\pi ix/N}.
\]
For $N=1$ this reads $\Lambda(1-s)=\Lambda(s)$, where $\Lambda(s)=\Gamma_{\mathbf R}(s)\zeta(s)$.
This is the functional equation of Riemann \cite{Rie59} and, for $N>1$, of
\cite[Chapter 9]{Dav00}.
\end{proposition}

\begin{proposition}\label{mt-zeta-functional-equation}
Let $s\in\mathbf C$ with $s\neq1$ and $s\notin\{0,-1,-2,\ldots\}$. Then
\[
 \zeta(1-s)=2(2\pi)^{-s}\,\Gamma(s)\cos(\pi s/2)\,\zeta(s),
\]
where $\zeta$ is the Riemann zeta function of Definition~\ref{mt-l-function}. This form of
Riemann's functional equation is in \cite[Chapter 2]{Tit86}.
\end{proposition}

\begin{proposition}\label{mt-strip-primitive-of-half-plane}
Suppose that $L(w,\eta)\neq0$ for every integer $M\geq1$, every Dirichlet character $\eta$ modulo
$M$ and every $w\in\mathbf C$ with $\operatorname{Re}w>1/2$ such that $\eta\neq\mathbf 1_M$ or
$w\neq1$ (Definitions~\ref{mt-dirichlet-character} and~\ref{mt-l-function}). Let $N\geq1$, let
$\chi$ be a primitive Dirichlet character modulo $N$, and let $s\in\mathbf C$ with
$0<\operatorname{Re}s<1$ and $L(s,\chi)=0$. Then $\operatorname{Re}s=1/2$.
\end{proposition}
\leanref{Formalism/GammaFactorNonvanishing/CritLine.lean}{Formalism.VW.critLine}
\begin{proof}
Since $0<\operatorname{Re}s<1$, we have $s\neq0$ and $s\neq1$. If $\operatorname{Re}s>1/2$, the
hypothesis for $\chi$ and $s$ contradicts $L(s,\chi)=0$.

Suppose that $\operatorname{Re}s<1/2$, and let $a=a_\chi$. The Gamma function has no zeros, and no
poles in the half-plane $\operatorname{Re}z>0$. So $\Gamma_{\mathbf R}(s+a)$ and
$\Gamma_{\mathbf R}(1-s+a)$ are finite and nonzero, and $\Lambda(s,\chi)=0$. The character
$\bar\chi$ is primitive modulo $N$, with $a_{\bar\chi}=a$, and its conjugate is $\chi$.
Proposition~\ref{mt-functional-equation} for $\bar\chi$ gives
\[
 \Lambda(1-s,\bar\chi)=N^{s-1/2}\,\varepsilon(\bar\chi)\,\Lambda(s,\chi)=0,
\]
so $L(1-s,\bar\chi)=0$. But $\operatorname{Re}(1-s)>1/2$ and $1-s\neq1$, so the hypothesis for
$\bar\chi$ and $1-s$ gives $L(1-s,\bar\chi)\neq0$.

Both cases are impossible, so $\operatorname{Re}s=1/2$.
\end{proof}

\begin{proposition}\label{mt-strip-of-half-plane}
Suppose that $L(w,\eta)\neq0$ for every integer $M\geq1$, every Dirichlet character $\eta$ modulo
$M$ and every $w\in\mathbf C$ with $\operatorname{Re}w>1/2$ such that $\eta\neq\mathbf 1_M$ or
$w\neq1$ (Definitions~\ref{mt-dirichlet-character} and~\ref{mt-l-function}). Let $N\geq1$, let
$\chi$ be a Dirichlet character modulo $N$, and let $s\in\mathbf C$ with $0<\operatorname{Re}s<1$
and $L(s,\chi)=0$. Then $\operatorname{Re}s=1/2$.
\end{proposition}
\leanref{Development/AutomorphicArithmetic/CharacterBoxes/ClaimStrip.lean}{Claim.strip}
\begin{proof}
Since $\operatorname{Re}s<1$, we have $s\neq1$. If $\operatorname{Re}s>1/2$, the hypothesis for
$\chi$ and $s$ contradicts $L(s,\chi)=0$.

Suppose that $\operatorname{Re}s<1/2$. Let $M$ be the conductor of $\chi$ and let $\psi=\chi^*$,
a primitive character modulo $M$ that induces $\chi$. Since $s\neq1$, Lemma~\ref{mt-euler-factors}
gives $L(s,\chi)=L(s,\psi)\prod_{p\mid N}(1-\psi(p)p^{-s})$. The left side is $0$, and each factor
of the product is nonzero by Lemma~\ref{mt-euler-factor-nonzero}, since $\operatorname{Re}s>0$. So
$L(s,\psi)=0$.

Let $a=a_\psi$. Since $\operatorname{Re}s>0$, the point $s$ is not a pole of $\Gamma_{\mathbf R}(s+a)$,
and $\Gamma_{\mathbf R}(s+a)\neq0$; as $s\neq0,1$, this gives
$\Lambda(s,\psi)=\Gamma_{\mathbf R}(s+a)L(s,\psi)=0$. The character $\bar\psi$ is primitive modulo
$M$, since $\bar\psi$ and $\psi$ have the same conductor, and $a_{\bar\psi}=a$ and
$\bar{\bar\psi}=\psi$. Proposition~\ref{mt-functional-equation} for $\bar\psi$ gives
\[
 \Lambda(1-s,\bar\psi)=M^{s-1/2}\,\varepsilon(\bar\psi)\,\Lambda(s,\psi)=0.
\]
Since $\operatorname{Re}(1-s)>0$, the factor $\Gamma_{\mathbf R}(1-s+a)$ is finite and nonzero, so
$L(1-s,\bar\psi)=0$. This contradicts the hypothesis for the character $\bar\psi$ modulo $M$ and the
point $1-s$, since $\operatorname{Re}(1-s)>1/2$ and $1-s\neq1$.

Both cases are impossible, so $\operatorname{Re}s=1/2$.
\end{proof}

\begin{proposition}\label{mt-riemann-of-half-plane}
Suppose that $\zeta(w)\neq0$ for every $w\in\mathbf C$ with $\operatorname{Re}w>1/2$ and $w\neq1$,
where $\zeta$ is the Riemann zeta function of Definition~\ref{mt-l-function}. Let $s\in\mathbf C$
with $\zeta(s)=0$. Suppose that $s\neq-2(n+1)$ for every integer $n\geq0$, and that $s\neq1$. Then
$\operatorname{Re}s=1/2$.
\end{proposition}
\leanref{Formalism/GammaFactorNonvanishing/RhOfRight.lean}{Formalism.VW.rh_of_right}
\leanref{Formalism/GammaFactorNonvanishing/CritLine.lean}{Formalism.VW.rh_of_every}
\begin{proof}
If $\operatorname{Re}s>1/2$, the hypothesis contradicts $\zeta(s)=0$, since $s\neq1$.

Suppose that $\operatorname{Re}s<1/2$. Then $s\neq0$, since $\zeta(0)=-1/2$. Put $w=1-s$. Then
$\operatorname{Re}w>1/2$, $w\neq1$ and $w\notin\{0,-1,-2,\ldots\}$, so
Proposition~\ref{mt-zeta-functional-equation} at $w$ gives
\[
 \zeta(s)=2(2\pi)^{-w}\,\Gamma(w)\cos(\pi w/2)\,\zeta(w).
\]
The factor $2(2\pi)^{-w}$ is nonzero, and the Gamma function has no zeros. If $\cos(\pi w/2)=0$,
then $w=2k+1$ for an integer $k$, so $s=-2k$. Here $k\leq-1$ contradicts
$\operatorname{Re}s<1/2$, $k=0$ contradicts $s\neq0$, and $k\geq1$ is excluded by hypothesis. So
$\cos(\pi w/2)\neq0$. Finally $\zeta(w)\neq0$ by hypothesis. Hence $\zeta(s)\neq0$, a
contradiction.

Both cases are impossible, so $\operatorname{Re}s=1/2$.
\end{proof}

\begin{remark}\label{mt-riemann-via-strip}
Suppose that $L(w,\eta)\neq0$ for every integer $M\geq1$, every Dirichlet character $\eta$ modulo
$M$ and every $w$ with $\operatorname{Re}w>1/2$ such that $\eta\neq\mathbf 1_M$ or $w\neq1$. The
conclusion of Theorem~\ref{mt-riemann-hypothesis} then also follows from
Proposition~\ref{mt-strip-of-half-plane} for $N=1$. Let $s$ be a zero of $\zeta$ with $s\neq1$ and
$s\notin\{-2,-4,\ldots\}$; then $s\neq0$, since $\zeta(0)=-1/2$. If $0<\operatorname{Re}s<1$,
Proposition~\ref{mt-strip-of-half-plane} for $\mathbf 1_1$ gives $\operatorname{Re}s=1/2$. The
function $\zeta$ has no zeros with $\operatorname{Re}s\geq1$, by Hadamard and de la Vallée Poussin
\cite{Had96,VP96}. If $\operatorname{Re}s\leq0$, then $s$ is not a pole of $\Gamma_{\mathbf R}$, so
$\Lambda(s)=\Gamma_{\mathbf R}(s)\zeta(s)=0$. Proposition~\ref{mt-functional-equation} for $N=1$
gives $\Lambda(1-s)=0$, and $\Gamma_{\mathbf R}(1-s)$ is finite because
$\operatorname{Re}(1-s)\geq1$. So $\zeta(1-s)=0$ with $\operatorname{Re}(1-s)\geq1$, which is
impossible.
\end{remark}
\leanref{Development/AutomorphicArithmetic/CharacterBoxes/ClaimStrip.lean}{Claim.zeta}
